\documentclass{article}
\title{論理学と証明 — 正しい推論の型を身につける}
\author{TeX 図書館}

\begin{document}

\section{この本の目的と読み方 — 推論の型を確かめる}

論理学とは、要するに\textbf{「推論の型を確かめる学問」}です。前提から結論を導き出すやり方には決まった型があり、その型が正しければ結論は必ず従います。

逆に型が壊れていれば、結論そのものは事実として正しくても、推論としては成立しません。まず 2 例から始めます。

\begin{quote}
\textbf{【例】}\ 「すべての犬は哺乳類である」「この動物は犬である」。第 2 文は第 1 文から得られるはずの結論のように見えますが、\textbf{推論としては成立していません}。
\end{quote}

第 1 文が保証するのは「この動物が犬なら哺乳類である」ことであって、「哺乳類だから犬である」とは言っていないからです。犬以外の哺乳類は無数にあり、きりんも馬も哺乳類です。

\begin{quote}
\textbf{【例】}\ ある薬の広告に「効果が確認されました」とあります。読んだ人が「この薬で治ったから、この薬が治した」と結論すれば\textbf{典型的な誤り}です。
\end{quote}

治った原因が薬とは限らず、自然に治った・検査の主観効果・別の薬を併用したことなど、他にも考えられるからです。

この 2 例に共通するのは、\textbf{前提が正しいのに、結論の出し方が型を外れている}という同じ型です。\textbf{結論の真偽と、推論の正しさは別の問題}です。

この 2 つを混ぜると、自分の主張を通すために論理を使ってしまいます。論理学が扱うのは、まさにその「型」のほうです。

\subsection{この本で使う記号は 5 つだけ}

道具は少ないほど使いやすく、最初から導入するのは次の 5 つだけです。

\begin{center}
\begin{tabular}{|c|c|l|}
\hline
記号 & 読み方 & 意味 \\
\hline
\(\Rightarrow\) & ならば・含意 & \(P\) が真ならば \(Q\) も真 \\
\(\Leftrightarrow\) & 同値 & 両方向が常に成り立つ \\
\(\land\) & かつ・連言 & 両方とも真 \\
\(\lor\) & または・選言 & 少なくとも片方が真 \\
\(\neg\) & ではない・否定 & 真偽を反転する \\
\hline
\end{tabular}
\end{center}

これに量化子 \(\forall\)(すべて)と \(\exists\)(存在する)が加わります。本書では真か偽かだけの古典論理を土台に、「未定」という 3 つ目を第 11 節で少しだけ触れます。

\begin{quote}
\textbf{【注意】}\ 右矢印の \(\Rightarrow\) が「ならば」で、両矢印の \(\Leftrightarrow\) が「同値」です。両者は区別します。「ならば」だけでは同値にはなりません。
\end{quote}

\subsection{読み方の進め方}

全 14 節の前半、第 1〜4 節で道具を立て、第 5〜8 節で証明の手法を学び、第 9〜14 節で応用の型とまとめを扱います。

方針は次の 3 つです。

\begin{enumerate}
\item \textbf{数式は最小限度}: 記号は 5 つと量化子 2 つだけ。計算ではなく\textbf{意味の読み方}に重点を置きます。
\item \textbf{具体例を先に}: 抽象論に入る前に、日常の文の訳を必ず置きます。
\item \textbf{誤りは必ず型で分かる}: 演繹の破れを典型的な誤謬として示します。
\end{enumerate}

記法は他の教科書と揃えています。集合と関数、\(\forall x\)、\(\exists x\) の書き方は『大学数学入門』第 1 節と同じものを使います。そちらで学んだ道具を、ここで呼び出して使うだけです。

\section{文の分解 — 述語と主語に切る}

論理学の入口は、普段すでに無意識に行っている作業を言葉にすることです。日本語の文は「主語と述語」という 2 つの部品に切れます。

すべての文をこの形にそろえれば、あとは機械的に扱えます。「すべての犬は哺乳類である」という文を 2 つに分けてみます。

\begin{itemize}
\item \textbf{主語}: すべての犬
\item \textbf{述語}: 哺乳類である
\end{itemize}

述語は、主語のうちのどの 1 つに「哺乳類である」という性質が結びつくかを表す規則です。これが本書における\textbf{関数}です。

『大学数学入門』第 1 節で扱った写像(規則)とまったく同じ扱いで、\textbf{「犬」という対象を「真か偽か」に写す規則}が述語関数です。

\begin{quote}
\textbf{【定義:述語関数】}\ 対象 \(x\) に対して「\(x\) が犬である」という性質を \(\mathrm{犬}(x)\) と書きます。\(\mathrm{犬}\) は 1 つの引数を取る関数で、値としてとるのは \textbf{真}か\textbf{偽}かの 2 つだけです。
\end{quote}

すべての対象を関数として扱うこと、これが述語論理の出発点です。

\subsection{「すべて」を書く}

「すべての犬は哺乳類である」における「すべての」が量化子 \(\forall\) の役目です。

「すべての対象の \(\mathrm{犬}\) が \(\mathrm{哺乳類}\) である」を記号で書くと次のようになります。

\[ \forall x \left( \mathrm{犬}(x) \Rightarrow \mathrm{哺乳類}(x) \right) \]

\textbf{見落としやすいのは、内側に \(\Rightarrow\) が入っている}ことです。「すべての犬が哺乳類である」は \(\mathrm{犬}(x) \land \mathrm{哺乳類}(x)\) ではありません。

犬でないものはこの文のFALSE 対象外だからです。含意の左辺が偽になれば、右辺がどうであれ含意は真になる、という第 3 節の定義がここで効きます。

\begin{quote}
\textbf{【注意】}\ 「すべての \(A\) は \(B\) である」を機械的に \(A(x) \land B(x)\) と書くのは典型的な誤りです。正解は \textbf{\(A(x) \Rightarrow B(x)\)} です。
\end{quote}

対比として「この犬は哺乳類である(この 1 匹だけについて)」なら \(\mathrm{犬}(x) \land \mathrm{哺乳類}(x)\) で正しいのです。同じ「犬は哺乳類」でも、量化詞の深さが違えば意味が変わります。

\subsection{「存在する」を書く}

「哺乳類が存在する」は \(\exists\) を使います。3 つの文を並べて置きます。

\begin{center}
\begin{tabular}{|l|l|}
\hline
日本語 & 記号 \\
\hline
すべての犬は哺乳類である & \(\forall x \left( \mathrm{犬}(x) \Rightarrow \mathrm{哺乳類}(x) \right)\) \\
ある哺乳類は青いである & \(\exists x \left( \mathrm{哺乳類}(x) \land \mathrm{青い}(x) \right)\) \\
この動物は犬ではない & \(\neg \mathrm{犬}(\mathrm{きりん})\) \\
\hline
\end{tabular}
\end{center}

\begin{quote}
\textbf{【例】}\ 「すべての犬は哺乳類である」を「すべての \(x\) は哺乳類である」と読むと誤りです。\(\forall x \left( \mathrm{哺乳類}(x) \right)\) は「哺乳類が存在する」という意味になり、きりんで反例になります。
\end{quote}

\textbf{量化詞の「すべて」が覆うのは述語の中まで}、という点が第 7 節の主題です。

\subsection{数式が壊れない書き方の作法}

表記にもきまりがあります。性質の日本語表記は \(\mathrm{犬}\) のように算用体で書き、\textbf{数式の中に数式だけを書きます}。句読点は数式の外に書きます。

日本語の関数名を数式の中に裸で書くと、数式のパーサが混乱するためです。次の 3 点を，守ります。

\begin{center}
\begin{tabular}{|c|c|l|}
\hline
意味 & 書き方(正) & 書き方(誤) \\
\hline
\(x\) が犬である & \(\mathrm{犬}(x)\) & 日本語の記号を 1 文字ずつ並べる \\
\(x\) と \(y\) の両方 & \(x \land y\) & 「および」の位置が曖昧 \\
必ず偽になる組合せ & \(P \land \neg P\) & \(P \lor \neg P\)(排中律) \\
\hline
\end{tabular}
\end{center}

第 4 節で表を使って演算を手入れし、第 5 節の証明へ進みます。そこまで来れば、「\(\sqrt{2}\) は無理数である」を自分で証明できます。

\section{命題論理の 4 演算 — 否定・連言・選言・含意}

命題とは、\textbf{真か偽のどちらかに決まる文}のことです。「すべての犬は哺乳類である」は量化子があるので厳密には述語論理の対象です。

第 2 節で量化子の内側だけを取り出した「\(\mathrm{犬}(x) \Rightarrow \mathrm{哺乳類}(x)\)」は 1 個の命題です。\textbf{この内側だけを取り出して扱うのが命題論理}で、4 つの演算を使います。

\subsection{否定 \(\neg\)}

\begin{quote}
\textbf{【定義:否定】}\ \(P\) の否定を \(\neg P\) と書きます。\(P\) が真なら \(\neg P\) は偽、\(P\) が偽なら \(\neg P\) は真です。日本語では「ではない」「〜でない」に当たります。
\end{quote}

否定を 2 回かけると元に戻ります。\textbf{二重否定}の法則です。

\[ \neg \neg P \equiv P \]

ここで \(\equiv\) は「同じもの」という意味です。この事実を真理値表で確認する方法は第 4 節で示します。

\subsection{連言 \(\land\) と 選言 \(\lor\)}

\begin{quote}
\textbf{【定義:連言と選言】}\ \(P \land Q\) は「\(P\) かつ \(Q\)」で、両方が真のときだけ真です。\(P \lor Q\) は「\(P\) または \(Q\)」で、少なくとも片方が真なら真です。
\end{quote}

\textbf{「または」は排他的ではありません}。両方が真でも \(P \lor Q\) は真です。日常の「A または B」は排他的に響きますが、論理では排他的とはしません。

日本語の接続語との対応は次のとおりです。

\begin{center}
\begin{tabular}{|c|c|l|}
\hline
日本語 & 記号 & 注意 \\
\hline
\(P\) かつ \(Q\) & \(P \land Q\) & 両方とも必要 \\
\(P\) または \(Q\) & \(P \lor Q\) & 片方だけでよい \\
\(P\) でも \(Q\) でも & \(P \lor Q\) & 「でも」は選言に変換 \\
\(P\) だが \(Q\) ではない & \(P \land \neg Q\) & 連言と否定の合成 \\
\(P\) も \(Q\) もでない & \(\neg (P \lor Q)\) & 二重否定に注意 \\
\hline
\end{tabular}
\end{center}

\subsection{含意 \(P \Rightarrow Q\) — 一番の難物}

含意には 2 つの意味があります。\textbf{「\(P\) ならば \(Q\)」}が日常の意味です。そして\textbf{「\(P\) が真なら \(Q\) も真」}が論理の定義です。この 2 つは一致しません。

\begin{quote}
\textbf{【定義:含意】}\ 命題の組 \((P, Q)\) に対し、\(P\) が真で \(Q\) が偽になる組合せが一つもないとき、\(P \Rightarrow Q\) は真です。\textbf{\(P\) が偽のときは、\(Q\) の真偽によらず含意は真}です。
\end{quote}

「\(P\) が偽なら何も言えない」のではなく、\textbf{「\(P\) が偽なら含意は真になる」}という点が、論理の混乱どころです。

日本語の「ならば」は日常では「\(P\) が真なら \(Q\) も真」という条件で使われますが、論理では真偽の組合せの定義として扱います。

「\(P\) が偽なら \(Q\) も偽になる」とは別のことで、それは \(\neg Q \Rightarrow \neg P\) という\textbf{対偶}です。第 5 節で扱います。

\begin{quote}
\textbf{【例】}\ \(P\) を「きりんである」、\(Q\) を「哺乳類である」とします。\(P\) は偽なので、\(Q\) の真偽によらず \(P \Rightarrow Q\) は真です。
\end{quote}

\textbf{前提が偽だと含意は常に真になる}。これが 2 値論理(古典論理)の重要な性質です。第 11 節で、\(P\) の値が「未定」の場合にもう 1 つの値を持つ 3 値論理を少しだけ紹介します。

\subsection{同値 \(\Leftrightarrow\)}

\begin{quote}
\textbf{【定義:同値】}\ \(P\) と \(Q\) の真偽が常に同じなら \(P \Leftrightarrow Q\) と書きます。「\(P\) ならば \(Q\) かつ \(Q\) ならば \(P\)」と言い換えられます。
\end{quote}

右矢印(片方向)と同値(両方向)は別物です。\textbf{片方向の「ならば」だけでは同値になりません}。これが最も多い失敗です。

\(P\) から \(Q\) への含意を証明しても、\(Q\) から \(P\) へ逆向き，含意を証明しなければ同値にはなりません。

\begin{quote}
\textbf{【注意】}\ 「ならば」\textbf{には理由(根拠)が要る}という感覚を覚えておくと、第 9 節の誤謬に効きます。「\(A\) ならば \(B\)」と「\(B\) ならば \(A\)」を同値だと信じるのは、論理的根拠が足りない典型的な誤りです。
\end{quote}

\begin{quote}
\textbf{【演習】}\ 次の日本語を記号に翻訳せよ。(1)「犬は動物である」 (2)「すべての犬は動物ではない」 (3)「犬も猫も哺乳類である」 (4)「犬は動物だが、猫は草食動物である」 (5)「もし雨が降るならば、湿度が高くなる」 (6)「東京は日本の首都ではない」
\end{quote}

\textbf{解答の指針:}\ (1) \(\forall x \left( \mathrm{犬}(x) \Rightarrow \mathrm{動物}(x) \right)\)、(2) \(\forall x \left( \mathrm{犬}(x) \Rightarrow \neg \mathrm{動物}(x) \right)\)、(3) \(\forall x \left( (\mathrm{犬}(x) \lor \mathrm{猫}(x)) \Rightarrow \mathrm{哺乳類}(x) \right)\)。

(4) \(\forall x \left( \mathrm{犬}(x) \land \mathrm{動物}(x) \right) \land \forall y \left( \mathrm{猫}(y) \Rightarrow \mathrm{草食動物}(y) \right)\)、(5) \(P \Rightarrow Q\)、(6) \(\neg \mathrm{首都}(\mathrm{東京})\)。(1) と (2) の違いが要点です。

\section{真理値表 — 16 個のマスを全部確かめる}

記号の演算は\textbf{表}で機械的に確かめられます。\(P\) と \(Q\) の真偽の組合せは 4 通りしかありません。

\(R\) も入れると \(2^3 = 8\) 通り、必要なら \(S\) まで入れて 16 通りになります。すべての演算結果を表の中に書き切れます。

\subsection{4 通り × 4 列の表}

\begin{center}
\begin{tabular}{|c|c|c|c|c|c|c|}
\hline
\(P\) & \(Q\) & \(\neg P\) & \(P \land Q\) & \(P \lor Q\) & \(P \Rightarrow Q\) & \(P \Leftrightarrow Q\) \\
\hline
真 & 真 & 偽 & 真 & 真 & 真 & 真 \\
真 & 偽 & 偽 & 偽 & 真 & 偽 & 偽 \\
偽 & 真 & 真 & 偽 & 真 & 真 & 偽 \\
偽 & 偽 & 真 & 偽 & 偽 & 真 & 真 \\
\hline
\end{tabular}
\end{center}

この表から次のことが一目で読めます。

\begin{itemize}
\item 「\(P\) が偽のとき \(P \Rightarrow Q\) が真」——3 行目と 4 行目
\item 「\(P\) が偽のとき \(P \Leftrightarrow Q\) は \(Q\) の真偽による」——3 行目は偽
\item 「\(P \lor \neg P\) は常に真」——\(P\) の列と \(\neg P\) の列のうち必ず片方が真
\end{itemize}

\subsection{対偶が同値であることの証明}

第 1 節で約束した、\textbf{対偶が同値}という事実を、この表で確かめます。「\(P \Rightarrow Q\) の対偶」は \(\neg Q \Rightarrow \neg P\) です。

\begin{center}
\begin{tabular}{|c|c|c|c|c|c|}
\hline
\(P\) & \(Q\) & \(P \Rightarrow Q\) & \(\neg Q\) & \(\neg P\) & \(\neg Q \Rightarrow \neg P\) \\
\hline
真 & 真 & 真 & 偽 & 偽 & 真 \\
真 & 偽 & 偽 & 真 & 偽 & 真 \\
偽 & 真 & 真 & 偽 & 真 & 真 \\
偽 & 偽 & 真 & 真 & 真 & 真 \\
\hline
\end{tabular}
\end{center}

3 列の「\(P \Rightarrow Q\)」と最後の「\(\neg Q \Rightarrow \neg P\)」が 4 行すべてで一致しています。よって

\[ P \Rightarrow Q \;\equiv\; \neg Q \Rightarrow \neg P \]

が\textbf{真理値表による証明}されました。行を全部並べて同じ値の並びになったもの同士は、常に同じ、という意味です。

\subsection{\(P \Rightarrow Q \equiv \neg P \lor Q\) の証明}

含意は「否定と選言」で書き換えられます。\(P \Rightarrow Q\) が偽になるのは「\(P\) が真で \(Q\) が偽」の 1 つだけです。

\begin{center}
\begin{tabular}{|c|c|c|c|c|c|}
\hline
\(P\) & \(Q\) & \(P \Rightarrow Q\) & \(\neg P\) & \(Q\) & \(\neg P \lor Q\) \\
\hline
真 & 真 & 真 & 偽 & 真 & 真 \\
真 & 偽 & 偽 & 偽 & 偽 & 偽 \\
偽 & 真 & 真 & 真 & 真 & 真 \\
偽 & 偽 & 真 & 真 & 偽 & 真 \\
\hline
\end{tabular}
\end{center}

\begin{quote}
\textbf{【定理:含意と選言の同値】}\ \(\;P \Rightarrow Q \equiv \neg P \lor Q\;\) である。これは真理値表の 4 つのマスがすべて一致することで証明される。この書き換えは第 5 節以降で頻出する。
\end{quote}

\subsection{4 つの法則}

\begin{quote}
\textbf{【定理:基本法則】}\ 次の 4 つは真理値表で一意に確かめられます。

\begin{enumerate}
\item \textbf{二重否定}: \(\neg \neg P \equiv P\)
\item \textbf{de Morgan の法則}: \(\neg (P \land Q) \equiv \neg P \lor \neg Q\)
\item \textbf{分配則}: \(P \land (Q \lor R) \equiv (P \land Q) \lor (P \land R)\)
\item \textbf{排中律}: \(P \lor \neg P \equiv \text{真}\)、\(\;P \land \neg P \equiv \text{偽}\)
\end{enumerate}
\end{quote}

de Morgan の法則を日本語に直すと、\textbf{「\(P\) かつ \(Q\) ではない」は「\(P\) ではない または \(Q\) ではない」}という意味です。「A も B もでない」を「A でない か B でない」に翻訳した形です。

\textbf{否定が 2 つあれば、1 つの否定にまとめられ、接続が反転する}。1 つめは連言から選言へ、2 つめは逆です。

\begin{quote}
\textbf{【注意:排中律の仮定】}\ \(P \lor \neg P\) が常に真なのは、\textbf{\(P\) が真でも偽でもないという状況を含んでいない}からです。3 値論理(第 11 節)ではこの排中律が揺らぐ。
\end{quote}

「\(P\) か \(\neg P\) のどちらかだ」という言い方も排中律の一部であり、証明の途中で仮定を無自覚に使うと破綻します。

\begin{quote}
\textbf{【演習】}\ 次の式を真理値表で調べ、\(P\), \(Q\) の値によらず常に同じになる組合せをすべて指摘せよ。(1) \(P \Rightarrow Q\) と \(\neg P \lor Q\) (2) \(P \land Q\) と \(Q \land P\) (3) \(\neg (P \lor Q)\) と \(\neg P \land \neg Q\) (4) \(P \Rightarrow Q\) と \(Q \Rightarrow P\) (5) \(\neg P \Rightarrow \neg Q\) と \(P \Rightarrow Q\)
\end{quote}

\textbf{解答の指針:}\ (1) 同じ(含意と選言の同値)。(2) 同じ(連言の交換則)。(3) 同じ(de Morgan)。(4) \textbf{違う}。\(P\) が真で \(Q\) が偽のとき、前者は偽、後者は真。(5) 同じ(\(P \Rightarrow Q\) の対偶)。

\section{反証法と対偶 — 「そうでない」ことを示す}

証明の方法は大きく 2 つに分かれます。\textbf{直接法}は、主張をそのまま順に示す方法です。\textbf{反証法}は、否定を仮定して矛盾を導く方法です。

数学の証明では、2 番目の反証法が最もよく使われます。理由は 1 つで、\textbf{「存在しない」を示すのは、「存在する」を示すよりずっと簡単}だからです。

存在するものを探すのは難しい場合があります。存在しないものを「1 つも無い」と示すのであり、割り切れない数を 1 つでも作れば終わるからです。

\subsection{対偶を使う}

第 4 節で証明した「\textbf{対偶は同値}」を使います。\(P \Rightarrow Q\) を証明したいときは

\[ \neg Q \Rightarrow \neg P \]

を証明すればよい、ということです。右辺の方がやりやすい場合があるため、実用上はこちらの手順で進めます。

\begin{quote}
\textbf{【定理:対偶で証明する】}\ \(P \Rightarrow Q\) を証明したいとき、代わりに \(\neg Q \Rightarrow \neg P\) を証明してよい。両者は真理値表で同じ並びになる(第 4 節)。
\end{quote}

たとえば「\(x^2 = 4\) ならば \(x = 2\) である」は同値ではありません。\(x = -2\) も \(x^2 = 4\) を満たすからです。対偶を使えば

「\(x \neq 2\) ならば \(x^2 \neq 4\) である」となります。\(x = -2\) は \(x^2 = 4\) を満たすので、矛盾が得られます。

\subsection{反証法を手順}

反証法は次の 4 手順です。

\begin{enumerate}
\item 証明したい命題を書き出す
\item その否定を仮定として導入する
\item 仮定と既存の事実から矛盾を導く
\item 矛盾が得られたので、否定は成立しないと結論する
\end{enumerate}

手順 2 と 3 が肝です。\textbf{仮定した「否定」が途中から忘れられてしまう}、という失敗が最も多いものです。

仮定を証明の途中で使わずにおしまい、その仮定を使わない議論をしてしまう、という展開が起きます。

\subsection{\(\sqrt{2}\) は無理数である}

ここから、実際に 1 つ証明します。代表性のある例なので、第 6 節の構成法と対比させて読むと理解が早いです。

証明したいのは次の命題です。

\[ \sqrt{2} \text{ は有理数ではない} \]

「無理数である」とは「有理数ではない」という意味です。有理数とは、2 つの整数 \(a, b\) を使って \(a/b\) と書ける数のことです( \(b \neq 0\) )。

\begin{quote}
\textbf{【定義:無理数と有理数】}\ 分母が \(1\) でない分数で表せる数を\textbf{有理数}、そうでない数を\textbf{無理数}という。平方根のうち、\(2\) のように整数でないものの平方根は必ず無理数になります。
\end{quote}

反証法で進めます。手順に従って書くと次のようになります。

\textbf{命題.} \(\sqrt{2}\) は有理数ではない。

\textbf{証明(反証法).} 逆に、\(\sqrt{2}\) が有理数であると仮定する。このとき、ある整数 \(a, b\) で \(b \neq 0\) が存在して

\[ \sqrt{2} = \frac{a}{b} \]

と書ける。両辺を 2 乗すると

\[ 2 = \frac{a^2}{b^2}, \quad \text{すなわち} \quad a^2 = 2b^2 \]

ここで \(b^2\) は 2 の倍数である。2 で割った余りを考えると、整数の 2 乗が 2 の倍数になるのは、元の整数が 2 の倍数のときだけだからである。

つまり \(a\) も 2 の倍数である。\(a = 2c\) と書くと、\(a^2 = 2b^2\) に代入して

\[ 4c^2 = 2b^2 \quad \Rightarrow \quad 2c^2 = b^2 \]

となる。すると \(b^2\) も 2 の倍数であるから、\(b\) も 2 の倍数である。つまり \(a\) も \(b\) も 2 の倍数なので、上の式を 2 で割ると

\[ \sqrt{2} = \frac{a}{b} = \frac{2c}{2d} = \frac{c}{d} \]

と、より分数の小さい形で書ける。すると同じ議論を繰り返すことで \(c, d\) も \(e, f\) も 2 の倍数になり、際限なく小さくなり続けなければならず、矛盾する。

ここで使ったのは、\textbf{整数には最小の 2 の倍数がある}という事実です。この矛盾から、仮定が誤りと分かります。

よって \(\sqrt{2}\) は有理数ではなく、\textbf{無理数である}。\(\square\)

\begin{quote}
\textbf{【注意:最少の因数分解】}\ 上の最後の議論は、2 ではなく任意の素数 \(p\) について同じです。「素数 \(p\) の倍数でない整数の \(p\) 乗は \(p\) の倍数にならない」という性質を使います。
\end{quote}

\subsection{そのほかの基本例}

同じ型が少し複雑な命題にも使えます。

「\(x^2 + y^2 = z^2\) ならば \(z > 0\) である」を対偶で示します。対偶は「\(z \le 0\) ならば \(x^2 + y^2 \neq z^2\)」です。\(z \le 0\) なら \(z^2 \ge 0\) で、一方 \(x^2 + y^2 \ge 0\) です。両方 0 になるなら \(x = y = z = 0\) ですが、これは通常的前提から除かれるので矛盾します。

\begin{quote}
\textbf{【演習】}\ 「素数 \(p\) の倍数でない整数の \(p\) 乗は、\(p\) の倍数ではない」を証明せよ。加えて、この性質を用いて「\(\sqrt{2}\) は無理数である」の証明を文章で説明せよ。
\end{quote}

\textbf{解答の指針:}\ \(a^2 \equiv 0 \pmod{2}\) と仮定して、\(a = 2k\) でも \(a = 2k+1\) でもなく、\(a \equiv 1 \pmod{2}\) から \(a^2 \equiv 1 \pmod{2}\) となることを示します。証明全体は「\(b^2 \equiv 0\) から \(b \equiv 0\)、\(b = 2d\) を代入、\(b^2 \equiv 0\) から再び \(d \equiv 0\)、矛盾」です。

\section{構成法と演繹 — 前提から結論を組み立てる}

第 5 節は「存在しないもの」を示しました。今度は逆に、\textbf{与えられた前提から結論を組み立てる}方法を見ます。「A かつ B」あるいは「A または B」の形から結論を引く、2 つの準則を使います。

\subsection{3 つの準則}

前提の形と結論の形、固定の 3 つです。

\begin{center}
\begin{tabular}{|c|c|l|}
\hline
前提 & 結論 & 呼び名 \\
\hline
\(A\) かつ \(B\) & \(A\) & 連言の除去 \\
\(A\) ならば \(B\) かつ \(A\) & \(B\) & 前件 modus ponens \\
\(A\) ならば \(B\) かつ \(\neg A\) & \(\neg B\) & 後件拒絶 \\
\hline
\end{tabular}
\end{center}

上の 2 番目と 3 番目が、第 8 節で扱う仮定法と否定後件の型そのものです。

演繹とは、\textbf{前提だけを手掛かりに結論を導く}ことです。結論そのものの真偽は問いません。前提が真なら結論も真になる、と保証できる推論を指します。

\begin{quote}
\textbf{【定義:演繹的妥当性】}\ 前提がすべて真なら結論も真になるとき、その推論は\textbf{妥当}( valid )であるという。結論が実際に正しいかとは無関係な、\textbf{形のの問題}である。
\end{quote}

\begin{quote}
\textbf{【例】}\ 前提「すべての犬は哺乳類である」「きりんは哺乳類である」から結論「きりんは犬である」を得る推論は、前提は真なのに結論が偽です。\textbf{妥当ではありません}。
\end{quote}

上の 3 つはいずれも\textbf{前提が真なら結論が必ず真}になります。真理値表で確かめれば、右の 3 つの型はいずれも例外なく成り立つのです。

\subsection{構成法の例}

構成法とは、\textbf{部品を 1 つずつそろえて、まとめて結論に持ち上げる}手法です。3 つの例で見ます。

\begin{quote}
\textbf{【例:連言から 1 つ取り出す】}\ 前提「\(A\) かつ \(B\)」から \(A\) だけを取り出すのは正当です。「きりんは哺乳類であり動物でもある」から「きりんは動物である」と結論するのは正しい型です。
\end{quote}

\begin{quote}
\textbf{【例:選言を分ける】}\ 前提「\(A\) または \(B\)」のとき、\(A\) が真なら \(B\) の真偽は不明ですが、\(\neg A\) が真なら \(B\) は真です。「雨が降っているか、風が吹いているか」で「雨は降っていない」と分かれば、風が吹いています。
\end{quote}

\begin{quote}
\textbf{【例:含意のつなぎ】}\ 前提「\(A \Rightarrow B\)」「\(B \Rightarrow C\)」から \(A \Rightarrow C\) が導けます。「\(A\) ならば \(B\)」「\(B\) ならば \(C\)」の 2 つから「\(A\) ならば \(C\)」です。
\end{quote}

\begin{quote}
\textbf{【注意】}\ 上の 3 番目を「\(A \Rightarrow C\) だけから \(A \Rightarrow B\) は導けない」と混同しないでください。含意のつなぎとは 2 つの含意を,\textbf{重ねて}使う操作です。
\end{quote}

\begin{quote}
\textbf{【演習】}\ 次の演繹の型を確かめ、結論が真になるような前提の例と、結論が偽になるような前提の例をそれぞれ 1 つ作れ。(1) 前提「\(P \Rightarrow Q\)」「\(Q \Rightarrow R\)」結論「\(P \Rightarrow R\)」 (2) 前提「\(P \land Q\)」結論「\(Q\)」
\end{quote}

\textbf{解答の指針:}\ (1) 前提が真なら必ず結論も真です。\(P\) を真、\(Q\) を真、\(R\) を真にすればよい。(2) 前提が真なら \(Q\) も真です。(3) \(\neg P\) が真なので \(P\) は偽、選言から \(Q\) が真になります。型が同じでも、前提の組合せで結論の真偽が変わる点が要点です。

\section{量化子 — 順序を交換できない}

第 2 節で \(\forall\) と \(\exists\) が出ました。重要なのは量化子の順序は、
そのままでは **交換できない** ことです。

\subsection{2 つの変数があるとき}

「すべての \(x\) と \(y\) について \(P(x, y)\) が成り立つ」は \(\forall x \forall y \, P(x, y)\) と書きます。

同じことを逆に並べた \(\forall y \forall x \, P(x, y)\) は、\textbf{意味が完全に同じ}です。両方とも「ありとあらゆる組の \((x, y)\)」を指すからです。

しかし \(\forall\) と \(\exists\) を混ぜると話は変わります。\(\forall y \exists x \, P(x, y)\) と \(\exists x \forall y \, P(x, y)\) は、\textbf{同じではありません}。

\begin{center}
\begin{tabular}{|l|l|}
\hline
記号 & 意味 \\
\hline
\(\forall y \exists x \, P(x, y)\) & どの \(y\) に対しても、それを満たす \(x\) が存在的(装置は \(y\) によってよい) \\
\(\exists x \forall y \, P(x, y)\) & ある 1 つの \(x\) が、どの \(y\) に対しても成り立つ \\
\hline
\end{tabular}
\end{center}

\subsection{反例で示す}

「\(P(x, y)\)」を「\(x > y\)」という関係に取ってみます。

自然数だけを考えると、\(\exists x \forall y \, (x > y)\) は偽です。「どの自然数より大きい自然数」は存在しないからです。

一方、\(\forall y \exists x \, (x > y)\) は真です。各 \(y\) に対して \(x = y + 1\) を選べばよいからです。選ぶ \(x\) が \(y\) ごとに変えてよいからです。

\begin{quote}
\textbf{【定理:量化子の順序は意味を変える】}\ \(\forall y \exists x \, P(x, y)\) と \(\exists x \forall y \, P(x, y)\) は一般に同値ではない。上の自然数の反例は証拠です。
\end{quote}

\subsection{日常の例}

同じ構造が日常にも出ます。「どの教室にも、机が 1 つ以上ある」は \(\forall y \exists x\) です。机の数は教室ごとに違うので、選ぶ机は教室ごとに変えてよいのです。

「この建物には、すべての部屋へ通じるエレベーターが 1 つある」は \(\exists x \forall y\) です。1 つの建物の中に、すべての部屋へ通じられるエレベーターが 1 つある、という意味になります。

前者は各教室に別々の机があっても成り立ち、後者は 1 つのエレベーターですべての部屋に行けることが要求されます。\textbf{順序が違うだけで要求されるものがまるで変わります}。

\begin{quote}
\textbf{【注意:否定と量化子】}\ 否定は量化子を裏返します。「すべての \(x\) が \(P\) ではない」は「\(P\) ではない \(x\) が存在する」の同値です。記号では \(\neg \forall x \, P(x) \equiv \exists x \, \neg P(x)\) と書きます。
\end{quote}

\begin{quote}
\textbf{【注意:存在の量化子を消さない】}\ 日本語の「すべての」を \(\forall x\) に訳すときは、\(x\) を実際に量化していることに注意します。「すべての犬は哺乳類である」を \(\mathrm{犬}(x) \Rightarrow \mathrm{哺乳類}(x)\) とだけ書くと、\(x\) が何なのかという量化の指定が欠けてしまいます。
\end{quote}

\begin{quote}
\textbf{【演習】}\ 次の日本語を \(\forall\),\(\exists\) を使って訳せ。(1)「すべての町に駅がある」 (2)「駅のある町が少なくとも 1 つある」 (3)「ある学生は、すべての科目で A の評価めている」 (4)「すべての科目の評価が A の学生が少なくとも 1 人いる」
\end{quote}

\textbf{解答の指針:}\ (1) \(\forall x \left( \mathrm{町}(x) \Rightarrow \mathrm{駅}(x) \right)\)、(2) \(\exists x \left( \mathrm{町}(x) \land \mathrm{駅}(x) \right)\)、(3) \(\exists x \left( \mathrm{学生}(x) \land \forall y \left( \mathrm{科目}(y) \Rightarrow \mathrm{評価A}(x, y) \right) \right)\)、(4) \(\exists x \forall y \left( \mathrm{科目}(y) \Rightarrow \mathrm{評価A}(x, y) \right)\)。(3) と (4) の差は、\(\forall\) が \(\exists\) の内側か外側かだけです。

\section{仮定法と否定後件 — 4 つの準則}

ここからは、含意を使った 2 つの準則を第 6 節でもう一度確認し、4 つにまとめます。\textbf{日常の議論の多くが、この 4 つの型で正当化できる}のが要点です。

\begin{center}
\begin{tabular}{|l|l|l|}
\hline
番号 & 形 & 呼び名 \\
\hline
1 & \(A \Rightarrow B\)、\(A\) / 結論 \(B\) & 仮定法 \\
2 & \(A \Rightarrow B\)、\(\neg A\) / 結論 \(\neg B\) & 否定後件 \\
3 & \(A \Rightarrow B\)、\(B\) / 結論 \(A\) & 肯定の誤り \\
4 & \(A \Rightarrow B\)、\(\neg B\) / 結論 \(\neg A\) & 否定の誤り \\
\hline
\end{tabular}
\end{center}

\textbf{1 と 2 は正しく、3 と 4 は誤り}です。真理値表で確かめると、1 と 2 は \(P\), \(Q\) のすべての値について結論が真になる一方、3 と 4 は \(P\) が真・\(Q\) が偽のとき結論が偽になります。

\subsection{仮定法の例}

「\(A \Rightarrow B\)」「\(A\)」から \(B\) 推出すのは仮定法です。日常では次のようになります。

\begin{quote}
\textbf{【例:仮定法】}\ 「この人がカードを持っていれば、入口に通される」「この人はカードを持っている」から「この人は入口に通される」。前提 2 が「カードを持っている」なので、含意の左辺が真です。
\end{quote}

\subsection{否定後件の例}

「\(A \Rightarrow B\)」「\(\neg A\)」から \(\neg B\) を導くのも正しい型です。日常では次のようになります。

\begin{quote}
\textbf{【例:否定後件】}\ 「切符がなければ入場できない」「この人は入場できた」から「この人は切符があった」。入場できたので「切符がない」は偽です。\textbf{含意の否定は両側の否定になる}、という形です。
\end{quote}

\begin{quote}
\textbf{【注意】}\ 「入場できない」の否定は「入場できる」です。日本語の二重否定を丁寧に取ると、この準則の正しさが分かります。
\end{quote}

\subsection{肯定の誤りと否定の誤り}

3 と 4 は\textbf{結論が真とは限りません}。同じ例で見ていくと分かります。

\begin{quote}
\textbf{【例:肯定の誤り】}\ 「カードがあれば通される」「この人は通された」から「この人はカードを持っていたことになる」。通されたという事実だけからは、カードを握っていたことは分かりません。別の経路があるからです。
\end{quote}

\begin{quote}
\textbf{【例:否定の誤り】}\ 「カードがなければ通されない」「この人はカードを持っていない」から「この人は通されない」。入場方法が 別にある場合、この結論は出てきません。含意の仮定の取り違えです。
\end{quote}

\begin{quote}
\textbf{【注意:逆含意の混同】}\ 「\(A \Rightarrow B\)」と「\(B \Rightarrow A\)」は一般に同値ではありません。「カードがあれば通される」の否定は「カードがなければ通されない」であり、「カードがなければ通されない」は「カードがあれば通される」と同値。これが否定後件が成り立つ理由です。
\end{quote}

\begin{quote}
\textbf{【演習】}\ 次の 4 つの推論のうち、どれが妥当か判定せよ。(1)「\(P \Rightarrow Q\)」「\(P\)」から \(Q\) (2)「\(P \Rightarrow Q\)」「\(Q\)」から \(P\) (3)「\(P \Rightarrow Q\)」「\(\neg P\)」から \(\neg Q\) (4)「\(P \Rightarrow Q\)」「\(\neg Q\)」から \(\neg P\)
\end{quote}

\textbf{解答の指針:}\ (1) 妥当(仮定法)。(2) \textbf{妥当ではない}(肯定の誤り)。\(P\) 偽・\(Q\) 真なら前提は真だが結論は偽。(3) 妥当(否定後件)。(4) \textbf{妥当ではない}(否定の誤り)。\(P\) 真・\(Q\) 偽なら前提は真だが結論は偽。4 つのうち 2 つだけが使える、というのが要点です。

\subsection{日常の例}

同じ 4 つの型が日常にもそのまま使えます。

\begin{quote}
\textbf{【例:交通】}\ 「赤信号なら止まれ」「この交差点の信号は赤ではない」から「この交差点を止まらない」。否定後件です。信号が赤ではないなら止まる必要はない、という結論になります。
\end{quote}

\begin{quote}
\textbf{【例:交通】}\ 「制限速度を少し超えただけでは摘発されない」「この車は摘発されていない」から「この車は制限速度を少し超えていない」。これは否定の誤りです。
\end{quote}

摘発されなかったのは、速度超過でなかったからとは限りません。基準をわずかに超えた場合や、ほかの違反が同時にあった場合にも摘発されないことがあります。

\begin{quote}
\textbf{【注意:結論が真でも推論は誤り】}\ 上の例では「速度超過していない」が事実だったとしても、推論の形自体が誤りであることに注目します。第 1 節の「哺乳類だから犬」と同じ考え方です。
\end{quote}

\section{論理学的誤謬 — 日常生活の中でよく見られる誤り}

ここからは、実際の議論に現れる誤りを見ます。誤謬とは、\textbf{推論の形そのものに破綻がある}ことです。結論が事実として正しいかどうかとは切り離して考えます。

\subsection{倒置論法と逆の誤り}

第 8 節で挙げた 2 つを、まず日本語で見直します。

\begin{quote}
\textbf{【誤謬:倒置論法(後件肯定)】}\  「\(A \Rightarrow B\)」「\(B\)」だから「\(A\)」とした。例:「テスト勉強した人は合格する」「あいつは合格した」「じゃあテスト勉強したんだな」。この以外の理由でも合格はします。
\end{quote}

\begin{quote}
\textbf{【誤謬:否定の誤り(前件否定)】}\  「\(A \Rightarrow B\)」「\(\neg A\)」だから「\(\neg B\)」とした。例:「雨が降れば湿度が上がる」「雨は降っていない」「なら湿度は上がっていない」。湿度上昇の原因は雨だけとは限らないからです。
\end{quote}

壊れているのは推論の形です。\textbf{「\(B\) が真だから \(A\) が真」は言えません}。第 8 節の表で見ると、\(A\) が偽で \(B\) が真のとき、前提は真であるのに結論は偽になります。

\subsection{偽の二者択一}

選択肢を 2 つに絞り、そのうち 1 つを否定することで、もう 1 つを強制する論法です。

\begin{quote}
\textbf{【誤謬:偽の二者択一】}\ 「この政策を批判するのは、無政府主義者だからだ」。政策を批判する立場が 2 つしかないと仮定しています。現実には「政策の一部だけを変える」という第 3 の立場もあります。
\end{quote}

論理的破綻は、\textbf{「\(A\) か \(B\)」の排他的な選言を、根拠なく仮定している}ことです。選択肢が 3 つ以上あるなら、この形は成立しません。

\subsection{人身攻撃}

相手を攻撃することで、その相手の主張を反証したつもりになる論法です。

\begin{quote}
\textbf{【誤謬:人身攻撃】}\ 「この案を主張しているのは、地方出身だ」「学問的な素養がない人の意見に反応するな」。話者の経歴と、主張の正しさは別の問いです。
\end{quote}

論理的破綻は、\textbf{「\(A\) が真であるかどうか」を「話者が \(A\) を主張したか」で代用している}ことです。話者を批判しても、主張の真偽は一歩も動きません。

\subsection{権威に訴える誤り}

権威を根拠にすること自体は、必ずしも誤りではありません。\textbf{その専門家が、その専門に限って話した主張}なら有力な根拠になります。

\begin{quote}
\textbf{【誤謬:権威に訴える誤り】}\ 「ある著名な学者がこう述べた。だから正しい」。発言した人の専門は、議題の分野ではない、という点が問題です。
\end{quote}

\textbf{必要なのは「その人の専門内の発言か」というチェックだけ}です。専門外の権威や、名前だけの権威を出しただけでは根拠になりません。

\subsection{傾き滑り論法}

小さな変更を一度認めてしまうと、最後には取り返しがつかないほど大きな変更まで進む、と主張する論法です。

\begin{quote}
\textbf{【誤謬:傾き滑り論法】}\ 「この小さな変更を認めたら、10 年後には取り返しがつかない大改革になる」。各段階で \(A \Rightarrow B\) の含意は真でも、\textbf{\(B \Rightarrow C\) が真である保証はない}。前提を 1 つ飛ばしています。
\end{quote}

論理的破綻は、\textbf{「\(A\) なら \(B\)」を「\(B\) なら \(C\)」に読み替えた}ことにあります。傾き滑りでは、\textbf{各段階の確からしさがあっても、累積した帰結まで確からしいわけではありません}。

\subsection{相関と因果の混同}

「\(A\) のとき \(B\) が起きる」を「\(A\) が \(B\) を起こす」と読み替える誤りです。

\begin{quote}
\textbf{【誤謬:因果の誤り(時間的前後を因果とみなす)】}\  「この薬を飲んだ人はたまたま治ったから、この薬が治した」。飲んだことと治ったことが同時に起きただけで、因果は示されていません。『統計学入門』第 2 節の例と同じ型です。
\end{quote}

\textbf{相関関係があることは、因果関係があることの証明にはなりません}。3 つめの要因が両方を動かす可能性があるからです。因果を示すには、実験や準実験の設計(『統計学入門』)が必要です。

\subsection{曖昧な言い回し}

曖昧な言葉を使って、都合のよい意味に取らせる手法です。

\begin{quote}
\textbf{【誤謬:曖昧化(はぐらかし)】}\  「まあ 当中 だろう」，「知っている人は知っているだろう」。何を指しているか決まっていないため、反論しようがありません。指示対象が確定していないことそのものが論理的破綻です。
\end{quote}

\subsection{そのほかの誤謬}

さらに次のような誤謬もあります。

\begin{itemize}
\item \textbf{言い換えによる反論}: 相手の主張を勝手に弱めてから反論する
\item \textbf{たどり着きの誤り}: 「その連鎖を追うと、最終的には取り返しがつかない」
\item \textbf{新しさへの訴求}: 新しいものほど良いとする
\item \textbf{分割の誤り}: 一群をすべて例外とみなして単純化
\end{itemize}

\begin{quote}
\textbf{【演習】}\ 次の 5 つの発言が、それぞれどの誤謬に当たるか指摘し、論理的な破れを説明せよ。(1)「この案に反対するのは、反対の側だけだ」 (2)「彼の意見が正しいのは、有名人だからだ」 (3)「アイスクリームが売れたから、アイスクリームが溺事故を起こす」 (4)「規制緩和を認めれば、最終的には国家が消滅する」 (5)「資料に書いてあるから正しい」
\end{quote}

\textbf{解答の指針:}\ (1) 偽の二者択一(第 3 の立場を無視)。(2) 権威に訴える誤り(専門の範囲を確認していない)。(3) 因果の誤り(第 3 の要因「夏」を無視)。(4) 傾き滑り論法(途中の含意を保証していない)。(5) 権威に訴える誤り(出典の内容を確認していない)。誤謬の名前よりも、\textbf{どの環が欠けているか}を言えることが本質です。

\section{証明の書き方 — 前提・仮定・結論の構造}

ここからは、実際の証明の書き方に入る。\textbf{証明とは「何が言いたいのか」を明確にして、仮定を一つずつ番号で置き、到達点まで推論の鎖で書くこと}です。

\subsection{3 つの部品}

証明には 3 つの部品があります。

\begin{itemize}
\item \textbf{前提}: 与えられた事実。証明の土台
\item \textbf{仮定}: 途中で置いて、証明の終わりに破棄するもの
\item \textbf{結論}: 言いたいこと
\end{itemize}

\subsection{前提に番号を付ける}

前提には番号を付けます。「第 1 行目の式は前提より、第 3 行目の式は仮定より」のように指せるようにするためです。

\begin{quote}
\textbf{【例:番号付きの証明】}\ \textbf{命題.} ある整数 \(n\) が偶数ならば、\(n^2\) は偶数である。\textbf{証明.} \(n\) が偶数であるとする。このとき、ある整数 \(k\) が存在して \(n = 2k\) と書ける。両辺を 2 乗すると
\end{quote}

\[ n^2 = (2k)^2 = 4k^2 = 2 \cdot (2k^2) \]

となる。(1) の定義より、\(n^2\) は \(2\) の倍数、すなわち偶数である。\textbf{証了}

\subsection{証明すべき命題の言い換え}

多くの証明は、\textbf{「何を示せばよいか」を見つける}ことから始まります。次の 4 つの型がよく使われます。

\begin{enumerate}
\item \textbf{定義を使う}: 「偶数である」の定義 \(n = 2k\) に照らし合わせる
\item \textbf{既知の定理を使う}: 既に証明した結果を前提として使う
\item \textbf{具体例で確かめる}: 代表例で型を確かめる
\item \textbf{反例を探す}: 否定を真にする例を探す(第 7 節の様式)
\end{enumerate}

\begin{quote}
\textbf{【注意:「 具体例で示す」は「全部で示す」ではない】}\ 1 つの例で一般の命題を証明できるのは、例外があるときだけです。偽の例が 1 つでも見つかれば、命題の否定が示せたことになります。
\end{quote}

\subsection{逆説でも同じ写法で示す}

第 5 節では「存在しないもの」を示しました。同じ論理を「存在するもの」に使うと、証明の型が反転します。

\begin{quote}
\textbf{【例:存在するものの証明】}\ 「\(n\) が大きいほど \(n^2\) も大きい」を示します。任意に \(x, y\) を取り、\(x > y\) とします。このとき
\end{quote}

\[ x^2 - y^2 = (x + y)(x - y) > 0 \]

となりますから \(x^2 > y^2\) です。任意に取った \(x, y\) について成り立つので、命題が証明できました。\textbf{証了}

\begin{quote}
\textbf{【注意:任意と存在の違い】}\ 「任意の \(x\) について \(P(x)\)」と「ある \(x\) について \(P(x)\)」は別の命題です。前者では \(x\) を好きな値に置き換えてよいので、\(x^2 - y^2\) の分解が使えます。後者では値が決まったままです。
\end{quote}

\begin{quote}
\textbf{【演習】}\ 「ある整数 \(n\) の 3 倍が 6 ならば、\(n\) は 2 である」を証明せよ。証明の型(直接法・対偶・反証法)のどれを使うかも述べよ。
\end{quote}

\textbf{解答の指針:}\ 3 番目の準則(否定後件)を使います。「\(3n = 6\) ならば \(n = 2\)」と「\(n \neq 2\)」から \(3n \neq 6\) が導かれることを示します。等式の割り算で矛盾するので、これが妥当な証明です。

\section{格子と束 — 集合の包含関係}

ここからは少し見点を変えて、\textbf{論理演算そのものを対象として}考えます。§3 節で使った 4 つの演算は、集合の演算にそのまま対応します。

\subsection{束と順序}

『大学数学入門』でも扱った概念を、ここでは論理の側から見直します。\textbf{束}とは、2 つの要素から「最小の共通の上位集合」と「最大の共通の上位集合」が一意に決まる構造です。

\begin{center}
\begin{tabular}{|c|c|l|}
\hline
論理 & 集合 & 意味 \\
\hline
\(A \land B\) & \(A \cap B\) & 共通部分(最小の共通の上位集合) \\
\(A \lor B\) & \(A \cup B\) & 和集合(最小の共通の上位集合) \\
\(\neg A\) & \(A^{c}\) & 補集合 \\
\(\top\) & \(\mathbb{U}\) & 全体の集合(常に真) \\
\(\bot\) & \(\varnothing\) & 空集合(常に偽) \\
\hline
\end{tabular}
\end{center}

この対応が成り立つ理由は、\textbf{含意が集合の包含関係に対応している}からです。「\(A\) ならば \(B\)」すなわち \(A \Rightarrow B\) は、\(A \subseteq B\) を意味します。

\subsection{3 値論理という第 3 の道}

ここまでは真と偽の 2 つだけを見てきました。けれども現実の推論には、「まだ分からない」という状況があります。

\begin{quote}
\textbf{【定義:3 値論理(strong Kleene 論理)】}\ 真・偽・未定の 3 つの値をとる論理です。否定は「未定」は「未定」のまま、「真」と「偽」を入れ替えます。連言は両方とも真なら真、両方とも偽なら偽、それ以外は未定です。
\end{quote}

\begin{quote}
\textbf{【注意:排中律が壊れる】}\ 3 値論理では、\(P\) が未定のとき \(P \lor \neg P\) も未定です。排中律が成り立たないため、\textbf{第 4 節で使った証明の一部が使えなくなる}ことがあります。「真か偽かどちらかだ」という前提を使う場面では、3 値論理を使うかどうかを明示する必要があります。
\end{quote}

\begin{quote}
\textbf{【演習】}\ strong Kleene 論理で、次の 3 つの値の組合せを調べよ。(1) \(P =\) 真、\(Q =\) 偽のときの \(P \Rightarrow Q\) (2) \(P =\) 未定、\(Q =\) 未定のときの \(P \land Q\) (3) \(P =\) 未定のときの \(\neg \neg P\)
\end{quote}

\textbf{解答の指針:}\ (1) 偽(左が真で右が偽なので含意は偽)。(2) 未定(両方が真でも偽でもないので未定)。(3) 未定(未定の否定は未定、その否定も未定)。3 つとも 2 値論理と同じ答えにはならない点が要点です。

\section{算術の論理 — 公理と証明}

論理学の応用として、\textbf{算術そのものを論理的対象にする}試みを簡単に紹介します。

\subsection{公理とは}

\textbf{公理}とは、\textbf{証明なしで受け入れることとした出発点}です。演繹の土台は、常に「これだけは仮定する」というまとまりとして置かれます。

\begin{quote}
\textbf{【定義:公理系】}\ ある理論の出発点として、指定された文の集まりです。\textbf{演繹の 6 節で見た仮定の集合にほかならない}ので、論理と数学の間に決定的な差はありません。
\end{quote}

\subsection{ペアノの公理}

自然数を「加算・乗算・順序の 3 つだけで的特征づける」最小の公理の集まりを、ペアノの公理と呼びます。

\begin{quote}
\textbf{【例:ペアノの公理の要点】}\ 1. \(0\) は数である。(2) 数があれば、次の数がある。(3) どの数も、次の数を持たない。(4) 加算の結合則、交換則が成り立つ。\(1\) は \(0\) の次、\(2\) は \(1\) の次、と定義する。
\end{quote}

この公理系の範囲で、\(1 + 1 = 2\) は定理として\textbf{証明できる} ます。「公理から有限回の推論で導かれる」という意味で、演繹的に導けます。

\subsection{ゴーデルの不完全性定理}

1931 年、ゴーデルは\textbf{「十分に強い公理系の中には、その公理系そのものを前提に誰も証明できない文が少なくとも 1 つ存在する}」という定理を示しました。

これは「\(1+1=2\) のような基本的な命題まで証明できない」という意味ではありません。\textbf{公理系そのものに関する文}が問題だからです。だからこそ「論理で何が証明できて、何ができない」のかを考える学問が面白いわけです。

\begin{quote}
\textbf{【注意:これは論理の限界ではない】}\ 不完全性定理は\textbf{論理自体の限界ではなく、公理系そのものの限界}を述べています。論理の規則そのものは健全です。
\end{quote}

\begin{quote}
\textbf{【演習】}\ 次の命題を論証し、必要なら第 5 節の反証法を使え。(1)「\(x\) と \(y\) の和が 5 ならば、\(x\) と \(y\) の和は奇数か偶数のどちらかである」 (2)「すべての素数は 2 または 3 で割り切れる」 ((2) は偽の命題である。反例を示せ。)
\end{quote}

\textbf{解答の指針:}\ (1) 5 自身は奇数。5 より小さい和なら、\(x\) が奇数か偶数のどちらかによって決まる。(2) 5 が素数だが 2 でも 3 でも割り切れないので反例になる。「すべての A は B」という命題は、\textbf{B でない A が 1 つでもあれば偽}であると気づけばよい、というのが要点です。

\section{論理学を使うとなぜ便利か — プログラミングとのつながり}

ここまでの道具は、\textbf{プログラムを書くときにもそのまま使えます}。関係は 3 つあります。

\subsection{条件分岐は含意そのもの}

プログラムの条件分岐は、そのまま含意の構造です。

\begin{lstlisting}[language=Python]
def can_enter(has_ticket: bool, is_staff: bool) -> bool:
    if has_ticket or is_staff:
        return True
    return False
\end{lstlisting}

このコードが保証しているのは「\(A \Rightarrow B\)」です。ここで \(A\) は「チケットがある、またはスタッフである」、\(B\) は「入場できる」です。

\begin{quote}
\textbf{【注意:関数の型は命題の型】}\ 「\(A \Rightarrow B\) ならば \(B\)」という定理は、コードでは「引数に \(A\) を受け取る関数は、必ず \(B\) を返す」という意味になります。\textbf{型注釈は命題の書き換え}に相当します。
\end{quote}

型検査と論理の関係が詳しく扱われているのは、『TypeScript 入門』です。ここでは「型検査とは、すべての実行経路で命題が成立するかを確かめること」という一点だけを取り上げます。

\begin{lstlisting}[language=Python]
def f(x: int) -> bool:
    # 型注釈は「x が int ならば、この関数は bool を返す」を意味する
    return x > 0

# 呼び出し側が期待する条件:
# f(3) が True であること
# f(-1) が False であること
\end{lstlisting}

\subsection{ループの不変条件}

ループの構造を論理的な型で表すと、検証しやすくなります。

\begin{lstlisting}[language=Python]
def max_sum(nums: list) -> int:
    best = nums[0]
    for v in nums:
        if v > best:
            best = v
    return best
\end{lstlisting}

ループの本体が走る条件は「\(v > \text{best}\) が真なら \(\text{best} = v\)」です。\textbf{この条件が常に真とは限らない}ことを明示しておくと、コードの検証がしやすくなります。

\begin{quote}
\textbf{【注意:テストは反例を探す作業】}\ テストで「この入力では true になる」を確かめるのは、\textbf{含意の左辺を真にして右辺を確かめる}ことです。バグを 1 つでも出せれば、命題の否定が示せます。
\end{quote}

\begin{quote}
\textbf{【演習】}\ 次の Python 関数の仕様を、含意と命題の形で書け。「与えられた整数 \(n\) が 10 以上ならば、関数は \(n - 10\) を返す」。また、この仕様を満たさない入力の例を 1 つ示せ。
\end{quote}

\textbf{解答の指針:}\(P\) を「\(n \ge 10\)」\(Q\) を「戻り値が \(n - 10\)」とすれば，含意「\(P \Rightarrow Q\)」が仕様です。\(n = 5\) を入れると \(P\) が偽なのに \(Q\) も偽になるので、この 1 例で「任意の \(n\) について成り立つ」ではなく「条件を満たすときだけ成り立つ」ことが分かります。

\section{総合演習と解答}

ここまでの道具をまとめて試します。\textbf{証明 3 問、誤謬の判別 4 問、命題論理の計算 5 問}の計 12 問です。まず自分で解いてから、解答の指針を開いてください。

\begin{quote}
\textbf{【総合演習:証明 3 問】}
\begin{enumerate}
\item 「すべての犬は哺乳類である」「この動物は犬である」から「この動物は哺乳類である」を、演繹的に示せ。
\item 「\(x\) と \(y\) は偶数ならば、\(x + y\) は偶数である」を証明せよ。
\item 「すべての素数は 2 または 3 で割り切れる」が偽であることを、反例を使って示せ。
\end{enumerate}
\end{quote}

\textbf{解答の指針:}\ 1. 「すべての犬は哺乳類である」を \(\forall x(\mathrm{犬}(x) \Rightarrow \mathrm{哺乳類}(x))\) と展開し、\(\mathrm{犬}(\mathrm{きりん})\) を代入して準則 2(仮定法)で \(\mathrm{哺乳類}(\mathrm{きりん})\) を導きます。2. \(x = 2m\), \(y = 2n\) と置いて \(x + y = 2(m+n)\) です。3. 5 は素数ですが 2 でも 3 でも割り切れないので反例です。\textbf{「すべての \(A\) は \(B\)」の否定は「\(B\) でない \(A\) がある」}ことを覚えておけば確実です。

\begin{quote}
\textbf{【総合演習:誤謬の判別 4 問】} 次の発言が、それぞれどの誤謬に当たるか指摘し、論理的に何が壊れているかを説明せよ。
\begin{enumerate}
\item 「検査で合格した人は勉強したそうだ」
\item 「この計画が通れば、将来は必ずしも悪くない」
\item 「あの専門家が推奨した方法だ」
\item 「この 2 つの事件は同じ夜に起こった」
\end{enumerate}
\end{quote}

\textbf{解答の指針:}\ 1. 倒置論法(後件肯定)。合格の原因は他にもある。2. 傾き滑り論法(途中から飛躍した)。3. 権威に訴える誤り(専門を確認していない)。4. 因果の誤り(時間的前後を因果とした)。\textbf{誤謬の名前より「どの環が欠けているか」が本質}です。

\begin{quote}
\textbf{【総合演習:命題論理の計算 5 問】} 次の 5 問を計算せよ。
\begin{enumerate}
\item \(P =\) 真、\(Q =\) 偽、\(R =\) 真のときの \(P \land (Q \lor R)\) の値
\item \(P \Rightarrow (Q \Rightarrow P)\) の真理値表を作り、常に真か調べよ
\item \(\neg(P \lor Q) \Leftrightarrow (\neg P \land \neg Q)\) が恒真か調べよ
\item \(P \lor (\neg P \land Q)\) を分配法則で変形せよ
\item 「\(P\) ならば \(Q\)」の否定を書き、de Morgan の法則で確認せよ
\end{enumerate}
\end{quote}

\textbf{解答の指針:}\ 1. \(Q \lor R\) は真、よって全体は真。2. 4 マスを埋めると常に真(恒真)。3. de Morgan の法則そのものなので恒真。4. \((P \lor \neg P) \land (P \lor Q) = P \lor Q\)。5. \(\neg(P \Rightarrow Q) \equiv P \land \neg Q\)。\(P \Rightarrow Q \equiv \neg P \lor Q\) の否定を取っても同じ結果になります。

\subsection{心構え}

最後に、この本を読んで 心得たいことを 1 つだけ挙げます。

\begin{quote}
\textbf{【心構え】}\ 論理学は人を疑わせる道具ではない。\textbf{推論の型を確かめ、敬意ある議論の共通言語を作るもの}である。論点を整理し、根拠を並べ、反論できる形で言えるようにするための道具であり、相手の動機を推測するための道具ではない。
\end{quote}

自分の主張を整え、\textbf{論証として}示すことを目的にします。誤謬を見いだすには、\textbf{まず相手の論の型を正確に書き出してから}欠けを指摘します。
\end{document}
